\documentclass[12pt]{article}
\usepackage[utf8]{inputenc}
\usepackage[T1]{fontenc}
\usepackage{hyperref}
\title{Antankarana Malagasy Language Resources for K-12 Curriculum Development}
\author{Compiled by Claude Sonnet 5.5 (claude-sonnet-5-5) for Andr\'as Kornai}
\date{}
\begin{document}
\maketitle

\section{Introduction}
Antankarana (Antakarana) Malagasy (ISO 639-3 xmv; Glottocode anta1255) is a northern Malagasy variety spoken in the Diana Region of Madagascar. Standard Malagasy, the official school language, is closely related to it. My searches found no dedicated dictionary, textbook or audio corpus for Antankarana. The resources below are general or comparative.

\section{Primary Educational Resources}
No Antankarana-specific primers, textbooks or school programmes were found. Standard Malagasy teaching resources, such as the University of Wisconsin LCTL page, could serve as a base for the closely related standard: \url{https://wisc.pb.unizin.org/lctlresources/part/malagasy/}

\section{Linguistic and Reference Materials}
\begin{itemize}
\item Wikipedia, ``Antakarana dialect'', including a vocabulary table compared with Standard Malagasy: \url{https://en.wikipedia.org/wiki/Antakarana_dialect}
\item Grambank entry: \url{https://grambank.clld.org/languages/anta1255}. It cites a study of the neighbouring Sakalava dialect, so the link to Antankarana is indirect.
\item ``Malagasy Dialects and the Peopling of Madagascar'': \url{https://arxiv.org/pdf/1102.2180}
\item ``The settlement of Madagascar: what dialects and languages can tell'': \url{https://arxiv.org/pdf/1107.4218}
\item Wikipedia, ``Malagasy language'': \url{https://en.wikipedia.org/wiki/Malagasy_language}
\end{itemize}

\section{Religious and Community Resources}
\begin{itemize}
\item Wikipedia, ``Antankarana'' (the people): \url{https://en.wikipedia.org/wiki/Antankarana}
\item A missionary-oriented profile of the Antakarana (to be read with that purpose in mind): \url{https://eu.aimint.org/unreached_people/antakarana/page/2}
\end{itemize}
No scripture translation specific to Antankarana was found.

\section{Digital and Audio Resources}
Language profile page: \url{https://detectlanguage.com/languages/xmv}. No audio or video was found.

\section{Curriculum Development Context}
Antankarana is spoken as a first language but is not known to be taught in schools. Documentation is thin, so a curriculum would need new community-based documentation alongside Standard Malagasy materials.
\end{document}